\documentclass{article}
\usepackage{graphicx}
\usepackage[margin=1.5cm]{geometry}
\usepackage{amsmath}

\begin{document}
\twocolumn

\title{Chapters 1-9 Warm-Up: Review}
\author{Prof. Jordan C. Hanson}

\maketitle

\section{Review Exercises}

\begin{enumerate}
\item Write a program that produces the following table of contents:
\begin{verbatim}
Python Programming: A First Course

Chapter 1: Getting Started....................1
Chapter 2: Variables and Data Types..........17
Chapter 3: Making Decisions..................35
Chapter 4: Repeating with Loops..............57
Chapter 5: Writing Functions.................81
\end{verbatim}
\item Write a program that collects information about a student:

\begin{verbatim}
score = int(input("Score: "))
sub = input("Submitted? yes/no: ") == "yes"
att = input("Attendance OK? yes/no: ") == "yes"
\end{verbatim}
Create the Boolean variable
\begin{verbatim}
ready = sub and att
\end{verbatim}
and use it in the following chained decision:
\begin{verbatim}
if score >= 90 and ready:
    print("Excellent")
elif score >= 70 and ready:
    print("Pass")
elif not sub or not att:
    print("Incomplete")
else:
    print("Revise")
\end{verbatim}
Run the program several times.  Find inputs that cause each of the four branches to execute.
\item Define the following \textit{diagonal} matrix:
\begin{verbatim}
A = [
    [1, 0, 0],
    [0, 1, 0],
    [0, 0, 1]
]
\end{verbatim}
Suppose we did not know the matrix was diagonal.  Create nested \verb+for+ loops and conditional logic that demonstrate the matrix is diagonal.  For example:
\begin{verbatim}
is_diagonal = True
for row in range(2):
    for column in range(3):
        if(...):
        		is_diagonal = False
        	...
\end{verbatim}
\item Define a function that evaluates $f(x) = \sin(x)^2 + \cos(x)^2$, and show that it is always 1, regardless of the input $x$.
\item Suppose a patient record is represented by a list:
\begin{verbatim}
patient = [1042,"Maria Garcia",37]
\end{verbatim}
The entries represent the patient ID, name, and age. Write a function
\begin{verbatim}
def display_patient(patient):
    # your code here
\end{verbatim}
that prints the record in the form
\begin{verbatim}
ID: 1042
Name: Maria Garcia
Age: 37
\end{verbatim}
\item Several course records can be stored inside one master list:
\begin{verbatim}
courses = [
 ["PHYS", 100, "SLC", 202],
 ["MATH", 150, "SLC", 105],
 ["BIOL", 101, "SLC", 210],
 ["PHYS", 200, "SLC", 202]
]
\end{verbatim}
Lists can be sorted using \texttt{sort()}. Define
\begin{verbatim}
def sort_courses(courses):
    courses.sort()
\end{verbatim}
Then run
\begin{verbatim}
sort_courses(courses)
for course in courses:
    print(course)
\end{verbatim}
\end{enumerate}

\end{document}